\textbf{(CAP-COLUNI 2018 -- Q20)} Considere a régua graduada com 5 centímetros de comprimento e as medidas
\[
\begin{aligned}
A&=8^{1/2}, & B&=\tan 30^\circ,\\
C&=(1+\sqrt3)^{-1}, & D&=\sqrt[3]{\sqrt{64}}.
\end{aligned}
\]
todas em centímetros.

Posicionando-se as medidas $A$, $B$, $C$ e $D$ na régua, é CORRETO afirmar que:

\includegraphics{./figs/2018-Q20.png}

\begin{enumerate}
	\item $A$ é irracional e está à direita de $B$, $C$ e $D$.
	\item $C+D$ é maior que $A$.
	\item $B$ é menor que $0{,}5$ e está entre $C$ e $A$.
	\item $D$ é menor que $B$ e é um número racional.
\end{enumerate}
